%%
%%  Annexes
%%
%%  Note: Ne pas modifier la ligne ci-dessous. / Do not modify the following line.
\ifthenelse{\equal{\Langue}{english}}{
	\addcontentsline{toc}{compteur}{APPENDICES}
}{
	\addcontentsline{toc}{compteur}{ANNEXES}
}
%%
%%
%%  Toutes les annexes doivent être inclues dans ce document
%%  les unes à la suite des autres.
%%  All annexes must be included in this document one after the other.
\Annexe{THERMODYNAMIC BALANCE EQUATIONS}
\label{app:thermodynamic_model}

This appendix provides the detailed expressions associated with the mass- and energy-balance equations introduced in Section~\ref{sec:review_thermodynamics}. The balances are applied locally to a surface control volume containing the ice layer and, when present, a liquid-water film.

\section{Mass Balance}

The local conservation of water mass is expressed as

\begin{equation}
    \dot{m}_{imp}
    + \sum \dot{m}_{in}
    =
    \sum \dot{m}_{out}
    + \dot{m}_{ev/su}
    + \dot{m}_{ice}
    \label{eq:app_mass_balance}
\end{equation}

where $\dot{m}_{imp}$ is the impinging-water mass flow rate, $\dot{m}_{in}$ and $\dot{m}_{out}$ are the incoming and outgoing runback-water flow rates, respectively, $\dot{m}_{ev/su}$ represents mass loss through evaporation or sublimation, and $\dot{m}_{ice}$ is the ice-accretion rate. The evaporation or sublimation rate may be evaluated as

\begin{equation}
    \dot{m}_{ev/su}
    =
    \frac{0.7}{C_{p,a}}\,HTC
    \left(
        \frac{p_{v,s}-p_{v,\infty}}{P_\infty}
    \right)
    \label{eq:app_evaporation}
\end{equation}

where $C_{p,a}$ is the specific heat capacity of air, $p_{v,s}$ and $p_{v,\infty}$ are the water-vapor pressures at the surface and in the freestream, respectively, and $P_\infty$ is the freestream static pressure. The vapor pressure at the water or ice surface is evaluated as~\cite{ozgen2009ice}

\begin{equation}
    p_v
    =
    3386
    \left(
        0.0039
        + 6.8096 \times 10^{-6}\,\bar{T}^{\,2}
        + 3.5579 \times 10^{-7}\,\bar{T}^{\,3}
    \right)
    \label{eq:app_vapor_pressure}
\end{equation}

where $p_v$ is expressed in Pa and

\begin{equation}
    \bar{T}
    =
    72 + 1.8\left(T-273.15\right)
    \label{eq:app_vapor_temperature}
\end{equation}

with $T$ expressed in K.

\section{Energy Balance}

Conservation of energy over the same surface control volume is written as

\begin{equation}
    \dot{Q}_{kin}
    + \dot{Q}_{d}
    + \dot{Q}_{a}
    + \sum \dot{Q}_{in}
    =
    \dot{Q}_{conv}
    + \dot{Q}_{rad}
    + \dot{Q}_{ev/su}
    + \sum \dot{Q}_{out}
    + \dot{Q}_{ice}
    \label{eq:app_energy_balance}
\end{equation}

where the individual contributions are detailed below. The kinetic energy associated with the impinging droplets is

\begin{equation}
    \dot{Q}_{kin}
    =
    \dot{m}_{imp}\frac{u_d^2}{2}
    \label{eq:app_qkin}
\end{equation}

where $u_d$ is the droplet impact velocity. The sensible heat associated with the impinging droplets is

\begin{equation}
    \dot{Q}_{d}
    =
    \dot{m}_{imp} C_{p,w}
    \left(T_s-T_\infty\right),
    \label{eq:app_qdroplet}
\end{equation}

where $C_{p,w}$ is the specific heat capacity of liquid water. The aerodynamic-heating contribution is expressed as

\begin{equation}
    \dot{Q}_{a}
    =
    \frac{r\,HTC\,U_\infty^2}
         {2C_{p,a}},
    \label{eq:app_qaero}
\end{equation}

where $r$ is the recovery factor. The convective heat exchange is

\begin{equation}
    \dot{Q}_{conv}
    =
    HTC\left(T_s-T_{rec}\right),
    \label{eq:app_qconv_recovery}
\end{equation}

The energy transported by incoming and outgoing runback water is

\begin{equation}
    \sum \dot{Q}_{in}
    =
    \sum \dot{m}_{in} C_{p,w}
    \left(T_s-T_{ref}\right),
    \label{eq:app_qin}
\end{equation}

and

\begin{equation}
    \sum \dot{Q}_{out}
    =
    \sum \dot{m}_{out} C_{p,w}
    \left(T_s-T_{ref}\right),
    \label{eq:app_qout}
\end{equation}

respectively, where $T_{ref}$ denotes the reference freezing temperature. Radiative heat transfer is given by

\begin{equation}
    \dot{Q}_{rad}
    =
    4\epsilon\sigma_r T_\infty^3
    \left(T_s-T_\infty\right),
    \label{eq:app_qrad}
\end{equation}

where $\epsilon$ is the surface emissivity and $\sigma$ is the Stefan--Boltzmann constant. The latent-energy contributions associated with evaporation and sublimation are

\begin{equation}
    \dot{Q}_{ev}
    =
    \dot{m}_{ev} L_v,
    \qquad
    \dot{Q}_{su}
    =
    \dot{m}_{su} L_s,
    \label{eq:app_qevsu}
\end{equation}

where $L_v$ and $L_s$ are the latent heats of vaporization and
sublimation, respectively. The energy associated with ice formation may be expressed as

\begin{equation}
    \dot{Q}_{ice}
    =
    \dot{m}_{ice}
    \left[
        C_{p,i}\left(T_s-T_{ref}\right)-L_f
    \right],
    \label{eq:app_qice}
\end{equation}

where $C_{p,i}$ is the specific heat capacity of ice and $L_f$ is the latent heat of fusion.
